%% annex/p07/calculi-and-cut-elimination-supplement.tex
% Necessity: the exact inherited proof sketches do not establish the distinct
% registered-calculus soundness and presentation-translation claims.
% Definition: the opening two items supply full editorial proofs of the
% registered-calculus soundness and presentation comparison. The later items
% prove LK--ND equivalence and canonical-LK completeness.
% Standing: it is separately attributed editorial evidence, not source proof
% text, a source-status upgrade, or a completed review determination.
% Nonredundancy: its body labels point only to the distinct registered-calculus
% additions and not to the exact source theorem or its proof sketch.
% Use: annex/p07.tex routes it once as editorial annex material, not restored
% source, with reciprocal body destinations for every item.

\section{Registered-calculus additions}

% Complete editorial proofs; the frozen source sketch is unchanged.
\begin{annexitem}{annex:p7.calculi.1}{thm:p07-lkd-soundness}{Soundness of the defined-negation sequent calculus}
\begin{proof}
Use exactly the rules of \hyperref[sheet:lk-defneg]{\texttt{sheet:lk-defneg}}.
A \emph{counterexample} to \(\Gamma\Rightarrow\Delta\) is a classical
structure \(\mathcal M\) and assignment \(s\) making every antecedent
occurrence true and every succedent occurrence false.  The latter condition
is vacuous for an empty succedent, as the theorem requires.  We prove by
induction on the finite derivation that its endsequent has no counterexample.
The induction hypothesis concerns all structures and assignments, so it may
be used after an assignment is varied in a quantifier case.

\emph{Initial sequents.} Identity \(\varphi\Rightarrow\varphi\) cannot
have a counterexample: its one formula would be both true and false.
For \(\bot L\), no assignment makes the antecedent occurrence of \(\bot\)
true, regardless of either context.

\emph{Weakening and contraction.} A counterexample to
\(\varphi,\Gamma\Rightarrow\Delta\) still makes all of \(\Gamma\) true
and all of \(\Delta\) false, contradicting the premise of WL.  A
counterexample to \(\Gamma\Rightarrow\Delta,\varphi\) similarly
contradicts the premise of WR.  Repeating a true antecedent occurrence
does not change its truth, and repeating a false succedent occurrence does
not change its falsity.  Thus a counterexample to a CL or CR conclusion is
also a counterexample to its premise.  Exchange needs no inference: both
contexts are multisets.

\emph{Cut and Mix.} Suppose
\(\Gamma,\Pi\Rightarrow\Delta,\Sigma\) has a counterexample.  Validity
of the first premise \(\Gamma\Rightarrow\Delta,\varphi^m\) forces
\(\varphi\) to be true, because every member of \(\Delta\) is false
and \(m\ge1\).  Hence all of \(\varphi^n,\Pi\) is true, while every
member of \(\Sigma\) is false, contradicting the second premise, since
\(n\ge1\).  This proves Mix soundness and, with \(m=n=1\), Cut
soundness.  It does not require the side contexts to omit \(\varphi\).

\emph{Conjunction.} For \(\land L\), a true
\(\varphi\land\psi\) makes both \(\varphi\) and \(\psi\) true, so a
counterexample to the conclusion refutes its premise.  For \(\land R\),
if the conclusion is false, all of \(\Gamma\) is true, all of \(\Delta\)
is false, and at least one of \(\varphi,\psi\) is false.  The premise
with that formula as its additional succedent is then false.

\emph{Disjunction.} For \(\lor L\), a counterexample makes
\(\varphi\lor\psi\) true; the true disjunct, together with the unchanged
contexts, gives a counterexample to the corresponding premise.  For
\(\lor R\), falsity of \(\varphi\lor\psi\) makes both displayed
succedent formulas of its premise false, so that premise is false.

\emph{Implication.} In \(\to L\), assume all of
\(\varphi\to\psi,\Gamma,\Pi\) is true and all of
\(\Delta,\Sigma\) is false.  The first premise
\(\Gamma\Rightarrow\Delta,\varphi\) forces \(\varphi\) true; truth
of the implication then forces \(\psi\) true.  This refutes the second
premise \(\psi,\Pi\Rightarrow\Sigma\).  For \(\to R\), falsity of
the principal implication makes \(\varphi\) true and \(\psi\) false,
giving a counterexample to
\(\varphi,\Gamma\Rightarrow\Delta,\psi\).  There are no primitive
negation cases on this sheet: \(\neg\varphi\) abbreviates
\(\varphi\to\bot\).

\emph{Term-instantiating quantifier rules.} The substitution theorem
\ref{thm:fol-substitution} gives, for every term \(t\),
\[
 \mathcal M,s\entails\varphi[t/x]
 \quad\Longleftrightarrow\quad
 \mathcal M,s[x\mapsto\sem{t}^{\mathcal M}_s]\entails\varphi.
\]
For \(\forall L\), truth of \(\forall x\,\varphi\) implies truth of
\(\varphi\) at this particular domain element, hence truth of
\(\varphi[t/x]\).  A counterexample to the conclusion therefore refutes
the premise.  For \(\exists R\), falsity of \(\exists x\,\varphi\)
implies falsity of \(\varphi\) at every domain element, including this
term value; the same equivalence makes \(\varphi[t/x]\) false and again
refutes the premise.  These arguments use the total capture-avoiding
substitution of the sheet, including any required bound-variable renaming.

\emph{Universal right.} Suppose the conclusion of \(\forall R\) is
false under \(\mathcal M,s\).  Choose \(a\in\dom{\mathcal M}\) such
that \(\varphi\) is false under \(s[x\mapsto a]\), and put
\(s_a=s[y\mapsto a]\), where \(y\) is the displayed eigenvariable.
Freshness and formula coincidence (Lemma~\ref{lem:fol-coincidence}) keep
every formula of \(\Gamma\) true and every formula of \(\Delta\)
false under \(s_a\).  Substitution identifies the value of
\(\varphi[y/x]\) under \(s_a\) with the value of \(\varphi\) under
\(s_a[x\mapsto a]\).  If \(y=x\), this last assignment is exactly
\(s[x\mapsto a]\).  If \(y\ne x\), freshness in the conclusion
\(\forall x\,\varphi\) implies \(y\notin\FV(\varphi)\), so these
two assignments agree on \(\FV(\varphi)\).  Coincidence therefore makes
the displayed premise instance false under \(s_a\), contrary to induction.

\emph{Existential left.} A counterexample to the conclusion makes
\(\exists x\,\varphi\) true.  Choose a witness \(a\) making
\(\varphi\) true under \(s[x\mapsto a]\), and again put
\(s_a=s[y\mapsto a]\).  The eigenvariable occurs free in neither
context, so coincidence preserves the antecedent-true and
succedent-false pattern.  If \(y=x\), the two updates coincide; otherwise
\(y\notin\FV(\varphi)\), by freshness in \(\exists x\,\varphi\).
Substitution and coincidence now make \(\varphi[y/x]\) true under
\(s_a\).  The premise consequently has a counterexample, again contrary
to induction.  This covers all rules, including empty contexts and the
eigenvariable case in which \(y=x\).
\end{proof}
\end{annexitem}

\begin{annexitem}{annex:p7.calculi.2}{thm:p07-lk-presentation-equivalence}{Canonical/defined-negation LK equivalence}
\begin{proof}
Use the exact calculi on Sheets~\ref{sheet:lk} and~\ref{sheet:lk-defneg}.
Multiplicities are
retained until an indicated contraction.  A formula of the defined-negation
language is also a canonical-LK formula, with no primitive negation symbol.

\emph{Syntactic facts about \(d\).} Formula induction gives
\[
 \operatorname{Var}(d\varphi)=\operatorname{Var}(\varphi),\qquad
 \FV(d\varphi)=\FV(\varphi),\qquad
 d(\varphi[t/x])=(d\varphi)[t/x].
 \tag{D}
\]
For the first two identities, atomic formulas and terms are unchanged;
the negation clause adds only \(\to\bot\), which adds no variable;
the other constructors preserve the same variable unions or remove the
same bound variable.  For the substitution identity, use strong induction
on formula complexity.  Atomic, falsity, and connective cases commute
directly with substitution.  In a quantified formula \(Qz\,\psi\),
both sides use the same substitution clause.  If it renames \(z\), the
least fresh variable is chosen outside the same set
\(\operatorname{Var}(\psi)\cup\operatorname{Var}(t)\cup\{x,z\}\),
by the first identity.  Apply induction first to that renaming and then
to substitution in the renamed matrix.  Its complexity is still that of
\(\psi\), by Lemma~\ref{lem:substitution-complexity}, so both uses of
induction are legitimate.  This proves literal equality for the fixed
deterministic substitution, rather than merely equality up to renaming.

\emph{Canonical LK to defined-negation LK.} Induct on the canonical
derivation, translating each endsequent occurrencewise.  This is also the
construction of Theorem~\ref{thm:p02-w02-p02-028-st-lk-translation}; its
rule cases are as follows.
\begin{enumerate}
\item Contexted identity is uncontexted identity
      \(d\varphi\Rightarrow d\varphi\), followed by one weakening for
      each occurrence in the two contexts.  Falsity left is the same
      initial sequent because \(d\bot=\bot\).
\item Left and right weakening and contraction use the corresponding
      rule on the translated occurrences.  Cut and Mix have the same
      split contexts on both sheets, so they apply to the translated cut
      or mix formula.  Further occurrences of that formula in the side
      contexts are permitted by both sheets.
\item For \(\neg L\), combine the translated premise
      \(d\Gamma\Rightarrow d\Delta,d\varphi\) with
      \(\bot\Rightarrow\), an instance of \(\bot L\), by split
      \(\to L\).  The result is
      \(d\varphi\to\bot,d\Gamma\Rightarrow d\Delta\), as required.
      For \(\neg R\), right-weaken the translated premise
      \(d\varphi,d\Gamma\Rightarrow d\Delta\) by \(\bot\), then
      use \(\to R\) to obtain
      \(d\Gamma\Rightarrow d\Delta,d\varphi\to\bot\).
\item The rules \(\land L,\land R,\lor L,\lor R,\to R\) have
      identical context arrangements on the two sheets.  Apply the rule
      of the same name; the required principal formula is its translation
      because \(d\) commutes with the connective.
\item For shared-context \(\to L\), the translated premises are
      \(d\Gamma\Rightarrow d\Delta,d\varphi\) and
      \(d\psi,d\Gamma\Rightarrow d\Delta\).  Split \(\to L\)
      yields
      \[
       d\varphi\to d\psi,d\Gamma,d\Gamma
       \Rightarrow d\Delta,d\Delta.
      \]
      Contract the duplicated antecedent and succedent occurrences to
      obtain exactly the translated conclusion.
\item For \(\forall L\) and \(\exists R\), (D) says that the term
      instance translates to \((d\varphi)[t/x]\). Apply the same rule
      with the same term. For
      \(\forall R\) and \(\exists L\), it identifies the instance
      with \((d\varphi)[y/x]\), and preservation of free variables
      in (D) preserves the eigenvariable condition on the entire
      conclusion.  Thus the same eigenvariable is legal.
\end{enumerate}
These cases exhaust the canonical sheet and prove the forward implication.

\emph{Embedding defined-negation derivations into canonical LK.}
Induct on a defined-negation derivation without changing any formula.
Its uncontexted identity is canonical identity with empty side contexts,
and its falsity initial sequent is canonical \(\bot L\).  The four
weakening/contraction rules, Cut, Mix,
\(\land L,\land R,\lor L,\lor R,\to R\) are unchanged instances.
So are \(\forall L,\exists R\), with the same total substitutions, and
\(\forall R,\exists L\), with the same eigenvariable condition and
unchanged conclusions.  It remains to simulate split \(\to L\).
Given its premises
\[
 \Gamma\Rightarrow\Delta,\varphi,
 \qquad \psi,\Pi\Rightarrow\Sigma,
\]
weaken the first by \(\Pi\) on the left and \(\Sigma\) on the right,
and the second by \(\Gamma\) on the left and \(\Delta\) on the right.
They become the two shared-context premises
\[
 \Gamma,\Pi\Rightarrow\Delta,\Sigma,\varphi,
 \qquad \psi,\Gamma,\Pi\Rightarrow\Delta,\Sigma.
\]
Canonical \(\to L\) now gives exactly
\(\varphi\to\psi,\Gamma,\Pi\Rightarrow\Delta,\Sigma\).
There are no primitive-negation rules to simulate in this direction.
In particular, embedding a derivation of
\(d\Gamma\Rightarrow d\Delta\) gives a canonical derivation of that
same sequent.  Recovering the original formulas requires one further step.

\emph{Formula bridges in canonical LK.} For every equality-free formula
\(\chi\), we construct derivations of both
\[
 \chi\Rightarrow d\chi
 \qquad\hbox{and}\qquad
 d\chi\Rightarrow\chi.
 \tag{B}
\]
Use strong induction on the complexity of \(\chi\), and, in the binary
cases, write \(\varphi'=d\varphi\), \(\psi'=d\psi\).
\begin{enumerate}
\item Atomic formulas and \(\bot\) are fixed by \(d\), so both
      sequents are identities.
\item For conjunction, weaken \(\varphi\Rightarrow\varphi'\) by
      \(\psi\), and \(\psi\Rightarrow\psi'\) by \(\varphi\), on
      the left.  Apply \(\land R\) and then \(\land L\), obtaining
      \(\varphi\land\psi\Rightarrow\varphi'\land\psi'\).
      The reverse construction uses the two reverse induction hypotheses.
\item For disjunction, right-weaken
      \(\varphi\Rightarrow\varphi'\) by \(\psi'\), then apply
      \(\lor R\).  Do the corresponding construction from
      \(\psi\Rightarrow\psi'\), weakening by \(\varphi'\).
      Apply \(\lor L\) to these two conclusions, obtaining
      \(\varphi\lor\psi\Rightarrow\varphi'\lor\psi'\).
      Interchanging primed and unprimed formulas gives the reverse sequent.
\item For implication, the induction hypotheses
      \(\varphi'\Rightarrow\varphi\) and
      \(\psi\Rightarrow\psi'\), with one right and one left
      weakening respectively, give
      \[
       \varphi'\Rightarrow\psi',\varphi,
       \qquad \varphi',\psi\Rightarrow\psi'.
      \]
      Shared-context \(\to L\) gives
      \(\varphi',\varphi\to\psi\Rightarrow\psi'\), and
      \(\to R\) gives
      \(\varphi\to\psi\Rightarrow\varphi'\to\psi'\).
      For the reverse sequent use
      \(\varphi\Rightarrow\varphi'\) and
      \(\psi'\Rightarrow\psi\) in the same construction.
\item For primitive negation, \(\varphi'\Rightarrow\varphi\)
      followed by \(\neg L\) gives
      \(\varphi',\neg\varphi\Rightarrow\).  Right weakening by
      \(\bot\) and then \(\to R\) give
      \(\neg\varphi\Rightarrow\varphi'\to\bot\).
      Conversely, combine \(\varphi\Rightarrow\varphi'\) with
      \(\varphi,\bot\Rightarrow\) by shared-context \(\to L\),
      obtaining \(\varphi,\varphi'\to\bot\Rightarrow\).
      Apply \(\neg R\) to obtain
      \(\varphi'\to\bot\Rightarrow\neg\varphi\).
\item For \(\forall x\,\varphi\), choose a variable \(y\) absent
      from \(\varphi\) and distinct from \(x\).  The formula
      \(\varphi[y/x]\) has smaller complexity than
      \(\forall x\,\varphi\).  Its two induction hypotheses, with
      (D), are
      \[
       \varphi[y/x]\Rightarrow\varphi'[y/x],
       \qquad \varphi'[y/x]\Rightarrow\varphi[y/x].
      \]
      Apply \(\forall L\) and then \(\forall R\) to the first,
      obtaining \(\forall x\,\varphi\Rightarrow\forall x\,\varphi'\).
      On the second, \(\forall L\) gives
      \(\forall x\,\varphi'\Rightarrow\varphi[y/x]\), and
      \(\forall R\) gives the reverse bridge.  Freshness follows from
      the choice of \(y\) and (D).
\item For \(\exists x\,\varphi\), use the same fresh \(y\) and
      the same two instance bridges.  Apply \(\exists R\) and then
      \(\exists L\) to the first to obtain
      \(\exists x\,\varphi\Rightarrow\exists x\,\varphi'\),
      and to the second to obtain the reverse bridge.  The conclusion
      sequents again contain no free \(y\), as required.
\end{enumerate}
This proves (B) for the whole equality-free grammar.

Finally embed the given defined-negation derivation to obtain
\(d\Gamma\Rightarrow d\Delta\) in canonical LK.  For each antecedent
occurrence \(\chi\) of \(\Gamma\), cut its bridge
\(\chi\Rightarrow d\chi\) against the corresponding occurrence in
that sequent.  This replaces exactly that occurrence of \(d\chi\) by
\(\chi\).  Then, for each succedent occurrence \(\chi\) of
\(\Delta\), cut against \(d\chi\Rightarrow\chi\).  The finite
sequence of Cuts yields \(\Gamma\Rightarrow\Delta\), with the original
multiplicities.  If either context is empty its replacement list is empty;
coinciding translations of different occurrences cause no problem, since
Cut selects an occurrence.  This proves the reverse implication.  These
constructions establish derivability equivalence; the final bridges use
Cut and establish no preservation of cut-freeness, height, or size.
\end{proof}
\end{annexitem}

\begin{annexitem}{annex:p7.calculi.3}{thm:p07-lk-nd-equivalence}{Derivability equivalence of canonical LK and classical natural deduction}
\begin{proof}
Write \(D_{\Delta}=\bigvee\Delta\), with the convention and fixed association
of Definition~\ref{def:p07-sequent-nd-translation}; in particular,
\(D_{\varnothing}=\bot\).  Repeated \(\lor I\) and \(\lor E\) give two finite
natural-deduction constructions that will be used below: if the support of
\(\Delta\) is contained in that of \(\Pi\), then
\(D_{\Delta}\vdash D_{\Pi}\), and a derivation of
\(D\lor A\lor B\) can be reassociated to a derivation of
\(D\lor(A\lor B)\).  When the smaller support is empty the first construction
is \(\bot E\).  These are derivations, not identifications of differently
associated formulas.

The sheet requires every discharge to be nonvacuous.  Whenever a construction
below has a derivation of \(C\) under which a displayed assumption \(A\) would
otherwise be unused, insert the finite detour
\[
 A\qquad C\quad\xRightarrow{\land I}\quad A\land C
 \quad\xRightarrow{\land E_2}\quad C.
 \tag{1}
\]
The marked occurrence of \(A\) then lies above the discharge.  Thus the usual
weakening admissibility is available without violating the exact nonvacuous
rule sheet.

First translate an LK derivation of \(\Gamma\Rightarrow\Delta\) into a
natural-deduction derivation of \(D_{\Delta}\) from
\(\operatorname{supp}(\Gamma)\), by induction on the LK derivation.  Identity
uses its displayed assumption followed by the finite-disjunction inclusion;
\(\bot L\) uses \(\bot E\).  Left weakening uses (1), right weakening uses
finite-disjunction inclusion, and either contraction changes only a
multiplicity forgotten by support.  For Cut, the translated left premise gives
\(D_{\Delta}\lor A\), and the translated right premise gives
\(D_{\Sigma}\) under assumption \(A\).  Eliminate the first disjunction:
its left branch injects \(D_{\Delta}\) into
\(D_{\Delta,\Sigma}\), while its right branch applies the second translation
and injects \(D_{\Sigma}\).  This derives the translated cut conclusion.
For Mix, with premises \(\Gamma\Rightarrow\Delta,A^{m}\) and
\(A^{n},\Sigma\Rightarrow\Pi\) (\(m,n\ge 1\)), the succedent
\(\Delta,A^{m}\) has the same support as \(\Delta,A\) and the antecedent
\(A^{n},\Sigma\) the same support as \(A,\Sigma\).  Both translations depend
on supports only, so the translated Mix premises are the translated
premises of the Cut on \(A\) with conclusion
\(\Gamma,\Sigma\Rightarrow\Delta,\Pi\), which is also the Mix conclusion; the
Cut construction then derives it.  This is the translation of the
contractions-then-Cut replacement of a Mix application.

The left logical rules translate by their matching eliminations.  For
\(\land L\), obtain both conjuncts and use the translated premise; for
\(\lor L\), apply \(\lor E\) to the two translated premises; for
\(\to L\), eliminate the translated formula \(D\lor A\): the \(D\)-branch
is already a required succedent, while the \(A\)-branch applies the open
assumption \(A\to B\), then the translated \(B\)-premise.  The rule
\(\neg L\) is the same case with contradiction and \(\bot E\).  The rules
\(\forall L\) uses \(\forall E\). For \(\exists L\), first translate
the premise by the induction hypothesis and then freshen that already
constructed finite ND subproof before applying \(\exists E\).
Choose a new variable absent from the whole translated proof, its final
context and its formulas; transpose the old eigenvariable with it,
carrying every term, formula and discharge occurrence coherently.
The explicit substitution covariance and literal alpha-restoration
construction in Annex~\ref{annex:p7.eqlk.nd} verifies these ordinary
logical steps without using equality. The LK proviso excludes only free
occurrences; this freshening supplies the stronger ND absence proviso.
No recursive translation call is made on the enlarged freshened proof.

For the right rules, \(\lor R\) is reassociation and injection, and
\(\exists R\) applies \(\exists I\) in the witness branch.  The translation of
\(\land R\) has premises \(D\lor A\) and \(D\lor B\).  Eliminate the first;
its \(D\)-branch is done, and in its \(A\)-branch eliminate the second, using
its \(D\)-branch directly and combining \(A,B\) by \(\land I\) in the other.
This yields \(D\lor(A\land B)\).

Three right-rule cases use classical RAA.  For \(\neg R\), the induction
hypothesis gives \(D\) under assumption \(A\).  Assume
\(\neg(D\lor\neg A)\).  From it derive \(\neg D\): an assumed \(D\) injects
into the forbidden disjunction.  Also derive \(\neg\neg A\): an assumed
\(\neg A\) injects into the same forbidden disjunction.  RAA gives \(A\), the
translated premise gives \(D\), and \(D\) contradicts \(\neg D\).  A final
RAA yields \(D\lor\neg A\).

For \(\to R\), the induction hypothesis gives \(D\lor B\) under assumption
\(A\).  Assume \(\neg(D\lor(A\to B))\) and derive \(\neg D\).  Under
assumption \(A\), eliminate \(D\lor B\); its \(D\)-branch gives
contradiction and then \(B\), and its other branch already gives \(B\).
Use (1) if needed and discharge \(A\), obtaining \(A\to B\), which contradicts
the outer assumption after injection.  RAA gives the required disjunction.

For \(\forall R\), the translated premise is \(D\lor A(y)\), where \(y\) is
free in neither the conclusion sequent nor its contexts.  Assume
\(\neg(D\lor\forall xA(x))\) and derive \(\neg D\).  Eliminating
\(D\lor A(y)\) gives \(A(y)\), since the \(D\)-branch yields contradiction
and then \(A(y)\).  The variable \(y\) is free in no remaining open assumption,
including the temporary negation, so \(\forall I\) gives \(\forall xA(x)\).
Injection contradicts the temporary assumption, and RAA completes the case.
The \(\neg R\), \(\to R\), and outer RAA assumptions are each used in the
displayed contradiction; hence their discharges are nonvacuous.  This exhausts
all rules of equality-free \texttt{sheet:lk}, including the four quantifier
rules, Mix and empty \(D\).

Conversely, induct on an equality-free derivation on
\hyperref[sheet:nd-fol]{\texttt{sheet:nd-fol}} to obtain a canonical-LK
derivation of \(\Gamma\Rightarrow A\) from every judgment
\(\Gamma\vdash A\).  The assumption case is contexted identity.  Contexts in
two- and three-premise natural-deduction rules are aligned by LK weakening and
their duplicate formula occurrences are removed by contraction.
The rules \(\land I\), \(\to I\), \(\neg I\), \(\forall I\), and
\(\exists I\) translate directly by the corresponding right rule (for
\(\neg I\), first cut the premise succedent \(\bot\) against
\(\bot\Rightarrow\), then use \(\neg R\)).  The eliminations
\(\land E_i\), \(\to E\), and \(\neg E\) cut the translated premises against,
respectively, the derived sequents
\[
 A\land B\Rightarrow A,\qquad
 A,A\to B\Rightarrow B,
 \qquad
 A,\neg A\Rightarrow\bot.
\]
The first is \(\land L\) over identity; the second is the \(\to L\) kernel
shown in Figure~\ref{fig:p07-sequent-nd-correspondence}; the third is
\(\neg L\) over contexted identity.  The rule \(\bot E\) cuts with the
initial sequent \(\bot\Rightarrow A\).

For \(\lor I_i\), weaken the succedent by the other disjunct and apply
\(\lor R\).  For \(\lor E\), align the two minor-premise contexts, apply
\(\lor L\) to obtain \(A\lor B,\Delta,\Sigma\Rightarrow C\), and cut the
translated major premise.  Universal elimination cuts against
\(\forall xA\Rightarrow A[t/x]\), obtained by \(\forall L\) over identity;
universal introduction is \(\forall R\), with the same freshness condition.
Existential introduction is \(\exists R\); existential elimination first
applies \(\exists L\) to the translated minor premise and then cuts the
translated major premise.  The eigenvariable condition is unchanged.
Alpha-conversion is admissible by induction on an LK derivation, consistently
renaming bound variables in every formula instance.  The two equality rules
are absent by the equality-free hypothesis.

For RAA, the translated premise is
\(\Gamma,\neg A\Rightarrow\bot\).  Cut it with \(\bot\Rightarrow\) to obtain
\(\Gamma,\neg A\Rightarrow\); apply \(\neg R\) to get
\(\Gamma\Rightarrow\neg\neg A\).  A cut-free derivation of
\(\neg\neg A\Rightarrow A\) is obtained by applying \(\neg R\) to identity
\(A\Rightarrow A\), producing \(\Rightarrow A,\neg A\), and then applying
\(\neg L\) to the displayed \(\neg A\).  Cut with it to obtain
\(\Gamma\Rightarrow A\).

The induction gives
\(\operatorname{supp}(\Gamma)\Rightarrow D_{\Delta}\) from a derivation of
\(D_{\Delta}\) in natural deduction.  Finally, canonical LK derives
\(D_{\Delta}\Rightarrow\Delta\) by induction on the finite support of
\(\Delta\): the empty case is \(\bot L\), the singleton case is contexted
identity, and the inductive case is \(\lor L\) applied to the two smaller
derivations; right weakening then restores any duplicate occurrences in
\(\Delta\).  Cut gives
\(\operatorname{supp}(\Gamma)\Rightarrow\Delta\), and left weakening restores
multiplicities in \(\Gamma\).  This proves the reverse implication and hence
the theorem.
\end{proof}
\end{annexitem}

\begin{annexitem}{annex:p7.calculi.4}{thm:p07-canonical-lk-completeness}{Finite-sequent completeness of canonical LK}
\begin{proof}
Soundness is exactly Theorems~\ref{thm:prop-lk-soundness} and
\ref{thm:fol-lk-soundness}.  For completeness, suppose first that the finite
propositional sequent \(\Gamma\Rightarrow\Delta\) is valid.  Then
\(\operatorname{supp}(\Gamma)\entails D_{\Delta}\), including the case
\(D_{\Delta}=\bot\).  The proved natural-deduction completeness theorem
\ref{thm:prop-completeness} gives
\(\operatorname{supp}(\Gamma)\vdash D_{\Delta}\), and the reverse translation
in Annex~\ref{annex:p7.calculi.3} gives the required LK derivation.

The first-order proof is the same.  The semantic condition on the sequent is
precisely
\(\operatorname{supp}(\Gamma)\entails D_{\Delta}\).  Under Choice,
Theorem~\ref{thm:fol-completeness} supplies a natural-deduction derivation; for
a finite or countable signature use
Corollary~\ref{cor:fol-completeness-countable}.  The equality-free
specialization of those complete Henkin proofs omits \(=I\), \(=E\), the
equality-congruence lemma, and the quotient of terms: the term structure uses
literal terms, while the connective, quantifier, consistency, witness, and
truth-lemma inductions are unchanged.  It therefore returns a derivation in
the equality-free restriction of \texttt{sheet:nd-fol}.  Annex~\ref{annex:p7.calculi.3}
translates that derivation to canonical LK.  The construction uses Cut in the
translation, so no cut-free strengthening follows.
\end{proof}
\end{annexitem}
